\scanpage{393}\mspage{132}

et qu'on veut qu'une formule (*) reste valable \inserted{pour tt $Y$ dans le domaine}, il faut ``en toute rigueur'' définir le $\mathrm{Hom\,ext}(a,b)$ non comme objet de $\mathrm{Hot}$, mais de $\mathrm{Ind\,Hot}$, \unsure{Il} se décrirait donc alors comme

\[
\mathrm{Hom\,ext}(a,b)=(\mathrm{Ch}_{n,\overline\Sigma}(a,b))_n
\quad\left|\begin{gathered}\text{Ind-objet}\\\text{dans }\mathrm{Hot},\end{gathered}\right.
\]

et l'objet $\mathrm{Ch}_{\infty,\overline\Sigma}(a,b)$ de $\mathrm{Hot}$ en est une sorte de limite inductive.

Je me suis demandé si la condition préalable, savoir que les \inserted{foncteurs} $H(a,b)$ \deleted{\unsure{fonc}} sur $\mathrm{Cat}$, transforment $\mathrm{Hot}$-équivalences $Y'\unsure{-|\longrightarrow} Y$ en bijections, était bel et bien une (forte) restriction. Je me suis convaincu que c'est bien le cas, déjà dans le cas particulier où $\overline\Sigma=\mathrm{Fl}(X)$. Il me semble que lorsque \inserted{la comp. connexe de} $\Omega(X,a)$ n'est \deleted{\unsure{belle}} pas contractile, la dite condition n'est pratiquement jamais satisfaite. J'ai pensé \inserted{p.ex.} au cas \deleted{d'une} \inserted{\unsure{où|d'où} \deleted{$\mathrm{Cat}$}} $X$ a le type \deleted{topologique} d'homotopie (dans $\mathrm{Hot}$) d'une $2$-sphère, de sorte que son $\Omega(X,a)=\mathrm{Ch}_\infty(X;a,a)$ doit avoir un $\pi_1\simeq\mathbb Z$. Prenons $Y$ ayant le type \unsure{dans|d'hos} $\mathrm{Hot}$ d'un cercle, disons

\diagram{\begin{tikzcd}[ampersand replacement=\&,column sep=small,row sep=small]
\&\&a\&\\
Y:\&\unsure{o|0}\arrow[ur]\arrow[dr]\&\&\unsure{o'|\bullet'|0'}\arrow[ul]\arrow[dl]\\
\&\&b\&\phantom{o},
\end{tikzcd}}

% archive sheet number: 392
